\documentclass[10pt]{article}
\usepackage[margin=.65in]{geometry}
\usepackage{amsmath,amssymb,hyperref}
\title{\vspace{-0.8cm}Certifying optimized transfer cutoffs}
\author{Collatz Lean research project}
\date{September 6, 2026}
\begin{document}
\maketitle\vspace{-0.4cm}
\begin{abstract}
We give a finite certificate for bounding every transfer request of the classical odd-run proof below a seed threshold. Two residue maxima per possible run length suffice. Lean verifies optimized cutoffs at parameters 61 and $2^{64}-1$, together with their maximality for this request criterion. All convergence conclusions retain the smaller-transfer hypotheses.
\end{abstract}
\section{Requests and the certificate}
Let $T$ be the shortcut Collatz map, $R(n)$ eventual arrival at one, and
$A(v)$ the implication $R(3v+2)\Rightarrow R(27v+20)$.
The bounded-transfer theorem needs only requests satisfying
\[
 a\ge2,\quad m>0\text{ odd},\quad 2^a m\le N,
 \qquad 36v+27=3^a m.                         \tag{1}
\]
If every such v is below U and $A(v)$ holds for all $v<U$, then every positive seed below N reaches one.

Choose L with $N<2^L$. Any request has $a<L$. For $r\in\{1,3\}$ define
\[
 M(B,r)=\begin{cases}r+4\lfloor(B-r)/4\rfloor,&r\le B,\\0,&r>B.\end{cases}
\]
The certificate checks, for $2\le a<L$ and each r,
\[
 3^a r\equiv3\pmod4\quad\Longrightarrow\quad
 3^a M(\lfloor N/2^a\rfloor,r)<36U+27.        \tag{2}
\]
This tests only the largest possible odd part in each eligible residue class. It does not enumerate all seeds or all transfer parameters.
\section{Soundness and descent}
For a request (1), let $r=m\bmod4$. Then r is 1 or 3 and $m\le\lfloor N/2^a\rfloor$. The residue-maximum formula gives $m\le M(\lfloor N/2^a\rfloor,r)$. Its congruence condition follows from (1). Equation (2) therefore yields $36v+27<36U+27$, hence $v<U$.

Compose this bound with the exact-request theorem in \texttt{BoundedTransfer.lean}. If a target $T^t(27U+20)$ lies below N, its positivity and the convergence cutoff prove $A(U)$ from the smaller-transfer hypotheses. This is a conditional induction rule, not an assertion that its hypotheses hold for all U.
\section{Checked cutoffs and maximality}
Lean's ordinary \texttt{decide} tactic verifies both certificates:
\[
 (U,N,L)=(61,415,9),\qquad
 (2^{64}-1,\ 30786325577727,\ 45).
\]
Their previous sufficient dyadic cutoffs were 128 and $2^{43}$ respectively.
For U=61, the next seed 415 has $a=5,m=13$ and requests $v=87\ge61$. For the larger U, seed 30786325577727 has $a=42,m=7$ and requests $v=21275914553349625401\ge U$.

Lean proves that every larger N, with any valid L, fails this request certificate in each case. Maximality concerns this arithmetic request system; it does not rule out other convergence proofs or larger cutoffs obtained by proving additional transfer instances.
The formal module \texttt{Collatz.RequestCutoff} uses only standard foundations. The Python optimizer proposes the checked outputs; its general correctness and a global descent or covering argument remain unproved. No literature-priority claim is made.
\end{document}
